\documentclass[10pt,conference]{IEEEtran}
\IEEEoverridecommandlockouts
% The preceding line is only needed to identify funding in the first footnote. If that is unneeded, please comment it out.

\usepackage{graphicx} % Required for inserting images
\usepackage{url}
\usepackage{listings}
\usepackage{xcolor}
\usepackage{amsmath}
\usepackage{amssymb}
\usepackage{flushend}
\usepackage{balance}


\graphicspath{{figs/}}

\usepackage{ifthen}
\newboolean{showcomments}
\setboolean{showcomments}{false}          
\ifthenelse{\boolean{showcomments}}
  {\newcommand{\nb}[2]{
    \fcolorbox{gray}{yellow}{\bfseries\scriptsize#1}
    {\scriptsize$\blacktriangleright${#2}$\blacktriangleleft$}
   } 
  }
  {\newcommand{\nb}[2]{}}
  
\newcommand{\jordi}[1]{\nb{Jordi}{#1}}
\newcommand{\javi}[1]{\nb{Javi}{#1}}
\newcommand{\adem}[1]{\nb{Adem}{#1}}
\newcommand{\gwendal}[1]{\nb{Gwendal}{#1}}

\abovedisplayshortskip=-10pt
\belowdisplayshortskip=-10pt
\abovecaptionskip=1pt
\belowcaptionskip=1pt

\floatsep=8.0pt plus 0pt minus 0pt
\textfloatsep=1.0pt plus 0pt minus 0pt
\intextsep=10.0pt plus 0pt minus 0pt

\makeatletter
\newcommand{\linebreakand}{%
  \end{@IEEEauthorhalign}
  \hfill\mbox{}\par
  \mbox{}\hfill\begin{@IEEEauthorhalign}
}
\makeatother

\title{Towards Automated Governance: A DSL for Human-Agent Collaboration in Software Projects%
\thanks{This work has been funded by the European Union under the Grant Agreement No 101189664 (MOSAICO project) and TED2021-130331B-I00 funded by MCIN/AEI/10.13039/501100011033 and European Union NextGenerationEU/PRTR. Views and opinions expressed are those of the author(s) only and do not necessarily reflect those of the European Union or the European Health and Digital Executive Agency (HADEA). Neither the European Union nor the granting authority can be held responsible for them.
Jordi Cabot is supported by the Luxembourg National Research Fund (FNR) PEARL program, grant agreement 16544475.}}



% Towards an Specification and enforcement of agent-aware governance policies: the case of Open Source Software
% \author{Anonymous (double blind)}
\date{June 2025}

\author{
    \IEEEauthorblockN{Adem Ait}
    \IEEEauthorblockA{\textit{University of Luxembourg}\\
    Esch-sur-Alzette, Luxembourg\\
    adem.ait@uni.lu}
    \and
    \IEEEauthorblockN{Gwendal Jouneaux}
    \IEEEauthorblockA{\textit{Luxembourg Institute of Science and Technology}\\
    Esch-sur-Alzette, Luxembourg \\
    gwendal.jouneaux@list.lu}
    \linebreakand    
    \IEEEauthorblockN{Javier Luis C\'anovas Izquierdo}
    \IEEEauthorblockA{\textit{Universitat Oberta de Catalunya} \\
    Barcelona, Spain \\
    jcanovasi@uoc.edu}
    \and 
    \IEEEauthorblockN{Jordi Cabot}
    \IEEEauthorblockA{\textit{University of Luxembourg} \\
    \textit{Luxembourg Institute of Science and Technology}\\
    Esch-sur-Alzette, Luxembourg \\
    jordi.cabot@list.lu}
}

\begin{document} 

\maketitle

\begin{abstract}
    The stakeholders involved in software development are becoming increasingly diverse, with both human contributors from varied backgrounds and AI-powered agents collaborating together in the process.
    This situation presents unique governance challenges, particularly in Open-Source Software (OSS) projects, where explicit policies are often lacking or unclear.
  %  While diversity can lead to more innovative solutions, it also requires robust governance mechanisms to ensure equitable representation and ethical decision-making.
    This paper presents the vision and foundational concepts for a novel Domain-Specific Language (DSL) designed to define and enforce rich governance policies in systems involving diverse stakeholders, including agents. 
    % Our DSL provides a structured approach for expressing governance rules, roles, responsibilities, and decision-making processes, with particular attention to fairness and representation. 
   % We propose core principles for such a DSL, emphasizing the explicit modeling of diverse participant profiles (both human and agent), their unique characteristics, roles, and adaptable decision-making processes, with a central focus on fostering fairness and transparency. 
    % By establishing a language tailored to modern governance needs, this work aims to enhance transparency, accountability, and inclusion in collaborative environments, fostering more equitable and effective interaction between humans and AI agents, while also providing a mechanism to automate the enforcement of these governance policies.
    % Our language supports modern governance needs by promoting transparency, accountability, and inclusion in collaborative environments.
    % It fosters equitable human-agent collaboration and enables automated enforcement of governance policies.
    This DSL offers a pathway towards more robust, adaptable, and ultimately automated governance, paving the way for more effective collaboration in software projects, especially OSS ones.
\end{abstract} 

\section{Introduction}
\label{sec:introduction}
%% Focus on diversity
% Prevalence of governance in the literature/OSS.
Governance is a key aspect of software development, especially in Open-Source Software (OSS), where the collaborative nature of the process requires clear guidelines and policies to ensure effective decision-making and accountability.
Governance policies help to establish roles, responsibilities, and processes for managing contributions, resolving conflicts, and maintaining the integrity of software projects~\cite{DBLP:journals/cacm/IzquierdoC23}.
Such policies make explicit the decision-making process, ensuring that all the involved stakeholders know how decisions are made and who is responsible for them.
However, most (OSS) projects do not explicitly define governance policies, and those that do are often poorly described and lack clarity~\cite{DBLP:journals/cacm/IzquierdoC23}.
Furthermore, when described, they are usually scattered across different project resources, making it difficult for contributors to understand the rules and procedures that govern the project~\cite{DBLP:conf/icse/IzquierdoC15}.

Recently, the landscape of software development is undergoing a profound transformation.
The rapid proliferation of AI-powered agents participating in development tasks, coupled with a growing recognition of the critical role of diverse human backgrounds, presents unprecedented challenges to these established governance paradigms~\cite{DBLP:journals/pacmhci/WesselSSWPCG18, DBLP:conf/xpu/RasheedWSKA0SA24, bjorn2023diversity, DBLP:journals/ese/Rodriguez-Perez21}.
AI agents, powered by advancements in Large Language Models (LLMs), are moving beyond simple automation to become active participants capable of complex communication, collaboration, and even decision-making~\cite{DBLP:conf/ijcai/GuoCWCPCW024}.
On the other hand, human diversity is also beneficial for software development, as it can lead to more innovative solutions and better decision-making~\cite{yang2022gender}, while also benefiting end-users~\cite{gunatilake2024impact}.
While this human-agent collaboration and broader human diversity promise significant benefits in innovation and productivity, they also expose a critical gap: the lack of governance frameworks designed to explicitly and holistically manage such multifaceted participation.\looseness-1

% Diverse world: humans, agents, minorities, etc.
% On the other hand, the software development landscape is becoming increasingly diverse.
% Recent research focuses on the diversity within the software engineering process~\cite{bjorn2023diversity,DBLP:journals/ese/Rodriguez-Perez21,DBLP:journals/corr/abs-2503-05470}.
% Diversity can be beneficial for software development, as it can lead to more innovative solutions and better decision-making~\cite{yang2022gender}, while also benefiting end-users~\cite{gunatilake2024impact}.
% Assessing this diversity in governance is crucial to ensure an equitable representation of all stakeholders.

% This diversity can be paired with the introduction of agents and bots that assist in the development process~\cite{DBLP:journals/pacmhci/WesselSSWPCG18, DBLP:conf/xpu/RasheedWSKA0SA24}.
% With the advancements of Large Language Models (LLMs), the agents can use the capabilities of these models to automate tasks, provide recommendations, and even participate in decision-making processes, thanks to their communication, collaboration and reasoning skills~\cite{DBLP:conf/ijcai/GuoCWCPCW024}.
% However, the adoption of these agents also raises concerns about the potential biases and ethical implications of their actions~\cite{DBLP:conf/models/MoralesCC24}.
% Appropriate identification of the agents' role in decision-making is one of the focuses of this work.

% With AI more fields are affected (not only software development).
% Furthermore, with the advancements in artificial intelligence (AI), particularly with the natural language comprehension of large language models (LLMs) and the increasing reliance on AI systems in various domains, the need for effective governance policies has become even more critical.
% The integration of AI agents into diverse fields, raise the concern for an holistic approach to governance that encompasses not only the technical aspects of software development but also the ethical, and social implications of AI systems.

% Motivation / proposal
% This paper first replicates the study of the presence of governance policies in OSS~\cite{DBLP:conf/icse/IzquierdoC15} to validate the need of explicit governance policies, and even more in the current agentic context. 
This paper first replicates a study on governance policies in OSS~\cite{DBLP:conf/icse/IzquierdoC15} to demonstrate the persistent need for explicit governance, particularly in today's agentic development context.
Then, we address this gap by proposing the conceptualization of a Domain-Specific Language (DSL) designed to define and enforce governance policies in systems involving a diverse set of stakeholders, covering both human and agent collaborators. 
% We illustrate the application of the DSL in the context of OSS development.
The proposed DSL provides a structured approach for expressing governance rules, roles, responsibilities, and decision-making processes.
% By focusing on the creation and definition of this language, we aim to enable system designers to clearly articulate governance requirements, ensuring that all stakeholders are represented while also agent interactions adhere to intended behaviors and constraints.
It encompasses various dimensions of governance, including decision-making procedures, positive and negative discrimination, and other methods to ensure a fair representation.
We envision our DSL to be field-agnostic, extensible, and adaptable to different contexts, allowing for the inclusion of various stakeholders and their specific needs.
We also provide tool support to automate and enforce the governance policies defined with our DSL.

% To adapt
The rest of the paper is structured as follows. 
Section \ref{sec:motivation} provides the motivation.
Section~\ref{sec:dsl} and \ref{sec:proof-of-concept}  presents the design and implementation of the DSL, and the proof of concept, respectively.
Section~\ref{sec:discussion} discusses the roadmap.
% Section N evaluates the DSL with a case study.
% Section~\ref{sec:related} provide the related work.
Finally, Section~\ref{sec:conclusion} concludes the paper and outlines the future work.

% \section{Background}
% \label{sec:background}
% \input{ase-background}

\section{State of the Art}
\label{sec:motivation}
% Part 1. 

\subsection{Diversity and Governance in Software Engineering}

Software engineering teams are becoming increasingly diverse across multiple dimensions.
Recent research has highlighted the importance of diversity within software development processes~\cite{bjorn2023diversity,DBLP:journals/ese/Rodriguez-Perez21}.
Diversity is being studied in terms of gender, race or ethnicity, but also in terms of cognitive diversity~\cite{DBLP:journals/corr/abs-2503-05470,DBLP:conf/icse/DuttaCST23}.
Studies have shown that diverse teams tend to produce more innovative solutions~\cite{DBLP:conf/icse/CatolinoPTSF19,doi:10.1073/pnas.2200841119}, and impact positively in software productivity~\cite{DBLP:conf/chi/VasilescuPRBSDF15}.

Adding to this human diversity, software projects now include a growing number of non-human participants in the form of bots and AI-powered agents~\cite{DBLP:journals/pacmhci/WesselSSWPCG18,DBLP:conf/xpu/RasheedWSKA0SA24}.
These automated participants have traditionally performed routine tasks like continuous integration, code quality checks, and dependency management~\cite{DBLP:conf/msr/KinsmanWGT21}.
% In the context of OSS, the use of automation mechanisms has become widespread, with GitHub Actions (GHAs) being adopted by approximately 30\% of the 5,000 most popular repositories~\cite{DBLP:journals/ese/WesselVGT23}.
% These automated tools have begun to play more significant roles in project governance by assisting with issue triaging, code review processes, and community management tasks such as welcoming new contributors and ensuring compliance with Contributor License Agreements~\cite{DBLP:conf/icsm/DecanMMG22}, or management of codes of conduct~\cite{DBLP:conf/icse/Cobos25}.
While these tools have traditionally handled routine tasks and basic governance functions~\cite{DBLP:journals/ese/WesselVGT23,DBLP:conf/icsm/DecanMMG22}, %, recent advances are enabling far more sophisticated capabilities.\looseness-1
% However, with the advancement of LLMs, agents have evolved beyond simple automation tools to become participants with advanced reasoning capabilities~\cite{xi2025rise}.
% LLM-based agents can now communicate, collaborate, and even participate in complex decision-making processes~\cite{DBLP:conf/ijcai/GuoCWCPCW024}.
% This evolution introduces a new dimension of diversity in software teams, the human-agent collaboration paradigm, which brings both opportunities and challenges to the governance.
the rise of LLMs is transforming these agents from mere automation tools into sophisticated collaborators capable of reasoning, communication, and participation in complex decision-making~\cite{xi2025rise, DBLP:conf/ijcai/GuoCWCPCW024}. 
This profound shift introduces an unprecedented dimension of diversity, presenting novel and urgent conceptual challenges for governance that existing frameworks are missing.

% While various proposals for AI governance have emerged from companies and countries in recent years~\cite{DBLP:journals/computer/ChestermanGHS24}, these often focus on AI systems as products rather than as participants in development processes.
% These proposals typically emphasize principles such as fairness, accountability, transparency, ethics, safety, and privacy.
% Existing AI governance frameworks, while valuable, predominantly address AI systems as products or services, focusing on principles like fairness, accountability, and transparency from an external or societal viewpoint~\cite{DBLP:journals/computer/ChestermanGHS24}. 
% However, they fall short in providing concrete mechanisms for integrating these principles into practical governance policies for diverse software teams that include both human and agent participants.
% In our proposal, we incorporate some identified indicators for AI:
% (1) autonomy level~\cite{DBLP:journals/nature/RahwanCOBBBCCCJ19}, to represent how independent the agent can make decisions;
% (2) explainability~\cite{DBLP:journals/inffus/ArrietaRSBTBGGM20}, usually understood for understandability, interpretability, explainability, and transparency of the AI system;
% and (3) human oversight~\cite{DBLP:journals/ijhci/Shneiderman20}, as whether the AI system is under human control or not.





% In the last few years, the number of proposals of governance policies in AI systems from companies and countries have quickly increased~\cite{DBLP:journals/computer/ChestermanGHS24}.
% % We could move the following to background (and maybe shorten it)
% From the analyzed documents, they identify six main principles: (1) fairness, which address biases, encompassed in three aspects: nondiscrimination, inclusiveness in impact, and inclusiveness in design; (2) accountability, seeking to minimize the risks associated with AI through the use of legal governance, risk management, and compliance mechanisms; (3) transparency and/or explainability, which focuses on facilitating human understanding of AI systems; (4) ethics/human-centricity, focuses on aligning AI with humans values by: benefiting humans, human control of technology, and human rights compliance; (5) safety/security, addresses the reliable operation of AI ensuring that AI systems operate without causing harm and by protecting from external threats; and (6) privacy, which ensures that data is properly handled, emphasizing human control over the data.
% In the context of our work, we focus on providing a DSL to specify governance policies for the decision-making, taking into account its diversity.
% Thus, we mainly focus on the principles of fairness, accountability, transparency and explainability. ELABORATE

\subsection{Governance in practice: the case of OSS development}
\label{sec:motivation:governance-oss}

A previous study illustrated the need for explicit governance policies in OSS~\cite{DBLP:conf/icse/IzquierdoC15} but lacked the mechanisms to address the challenges introduced by participant diversity, particularly the inclusion of AI agents in decision-making processes.
Our current proposal aims to fill this gap by incorporating diversity-aware governance mechanisms that acknowledge the different characteristics of participants and ensure equitable representation in decision-making processes.


% 2. Explicit definition of governance policies in top25 GH
While the landscape of participants is evolving, the explicitness and nature of governance policies in practice, particularly in OSS, remain a concern. 
To understand this current state, we replicated the study conducted by C\'anovas Izquierdo and Cabot
~\cite{DBLP:journals/cacm/IzquierdoC23}, which analyzed the top 25 starred software projects and found that most OSS projects do not have explicit governance policies, and those that do are often poorly defined and lack clarity.
The governance evidence is analyzed from the explicit indications of four dimensions: 
(1) whether there is some workflow to contribute, different from the typical pull-based development process; 
(2) who makes the decisions to accept code, and how; 
(3) how long it takes to review or accept a contribution; 
and (4) how to become a contributor.\looseness-1

% \begin{figure}
%     \centering
%     \includegraphics[width=0.8\columnwidth]{governance-top25.png}
%     \caption{Governance evidence in the top 25 \textsc{GitHub} projects.\adem{Could we report the results in text rather than in a figure?}}
%     \label{fig:governance}
% \end{figure}

% Figure~\ref{fig:governance} shows the results of our analysis.
We found that 68\% of repositories reported at least one of the four dimensions, 24\% reported none, and 8\% reported all four dimensions. %\footnote{Results available at: \url{https://doi.org/10.5281/zenodo.15856633}}.
The situation is showing slight improvements from the previous study.
In comparison, we have seen an increase in the number of projects that partially discussed governance, from 32\% to 68\%, mainly because of the report of the contribution process.
However, the span to review or accept a contribution is only defined in 16\% of the projects, how to become a contributor is only defined in 20\% of the projects, and who makes the decisions, and how, to accept code is only defined in 24\% of the projects.
Only 2 projects (8\%) provided a full description of the policies governing them. % mentioned all governance dimensions. 
Note that no project included a governance.md file but reported them in the contributing.md file or on their documentation website.
None of the analyzed projects adopted a DSL to specify their governance policies, which could help in the definition, enforcement, and automation of the governance policies~\cite{DBLP:journals/cacm/IzquierdoC23}.
% We can also say that we do not know whether these policies are being applied. We could use our DSL to test this in the Test Case Section.


\section{Defining Governance Policies}
\label{sec:dsl}
% To the best of our knowledge, there is no DSL for defining governance policies in systems that address the diversity of participants.
% Our contribution is a DSL to facilitate the definition of governance policies, while proposing constructs for specifying the profile of participants and their role.

% DSL in governance (motivation)
% In the context of governance, a DSL can serve as a formal framework for specifying policies, roles, and decision-making procedures, thereby enhancing the transparency and reliability of the interactions.
A DSL can formalize governance policies, roles, and decision-making procedures, enhancing transparency and reliability.
%By using a DSL, the developer can use domain-specific constructs and therefore address the problem more efficiently~\cite{DBLP:books/daglib/0034522}, which can be further automated and enforced.
% Furthermore, by implementing an engine to execute such DSL, the governance policies can be automatically enforced, reducing the risk of human error and ensuring that the systems operate within the defined constraints.
To the best of our knowledge, there is no DSL for defining governance policies in systems that address the diversity of participants.
%With our approach, governance policies can be defined in a more structured way, allowing for the specification of the participants involved, the scope of the policy, and the conditions that must be fulfilled. 
%This is particularly interesting in systems that include agents, as it allows for a clearer understanding of how these agents fit into the overall governance framework.

Several studies % have been conducted to 
analyze the governance of software projects~\cite{DBLP:conf/icse/ChulaniWY08,DBLP:conf/sigsoft/Herbsleb16}, particularly % in the context of 
OSS projects~\cite{DBLP:journals/ism/PeltJBO21,DBLP:conf/icse/IzquierdoC15,DBLP:conf/ease/KeertipatiLS16}.
However, to the best of our knowledge, only C\'anovas Izquierdo and Cabot~\cite{DBLP:conf/icse/IzquierdoC15} proposed a DSL, where they covered the main dimensions  %they identified 
to define and enforce governance rules in OSS projects.
Note, their proposal did not provide support to assess the diversity of the participants involved in the decision-making process.
% Our DSL extends their proposal by providing a more expressive syntax to define the profile of participants involved in the decision-making process and their impact on the decision.
Our proposal completely revamps that DSL and adapts it to the current reality of software development involving participants with different profiles, the possibility of having uncertain AI agents, and much more complex decision procedures with different participant weighting systems.

To collect the requirements for our DSL, we have carefully considered existing AI governance frameworks and principles~\cite{DBLP:journals/computer/ChestermanGHS24}. 
%While these frameworks provide valuable guidance on fairness, accountability, and transparency, our DSL operationalizes these principles for practical application in diverse software teams that include both human and agent participants. 
For instance, we incorporate key AI governance indicators identified in the literature, such as autonomy level~\cite{DBLP:journals/nature/RahwanCOBBBCCCJ19} and explainability~\cite{DBLP:journals/inffus/ArrietaRSBTBGGM20}%
%, and human oversight~\cite{DBLP:journals/ijhci/Shneiderman20}
, thus integrating them directly into our language constructs for agent representation.

% DSLs are languages specially designed to help to solve a problem in a particular domain.
Based on these frameworks and our own experience in open-source development and the analysis of open-source communities, we propose, in what follows, our governance DSL.
A DSL is composed of three main elements~\cite{DBLP:books/sp/WasowskiB23}: 
(1) abstract syntax, which defines the concepts and relationships of the domain where the language is applied; 
(2) concrete syntax, which defines the notation of the language (e.g., textual, diagram-based, etc.); and 
(3) semantics, which defines the meaning of the language constructs.
% Furthermore, DSLs can be classified into external DSLs, which are generally defined by a grammar; and internal DSLs, which are embedded in a general-purpose programming language (known as host language).
In the following, we discuss each component. 

\subsection{Abstract Syntax}
\label{sec:our-dsl:abstract-syntax}

The abstract syntax of our DSL is defined via the metamodel shown in Figure~\ref{fig:metamodel}. 
Metamodels restrict the structure of valid models (i.e., DSL instances, in our case, concrete policies) and define the relationships between the different elements of the language (i.e., the different concepts that can be used to define a policy and how these concepts can be linked to each other)~\cite{DBLP:series/synthesis/2017Brambilla}.
%We have defined a governance policy metamodel that captures the key concepts and relationships involved in defining governance policies.
%The metamodel is designed to be simple and intuitive, allowing for easy understanding and use by practitioners.
Governance models conforming to this metamodel represent specific sets of governance policies.

\begin{figure*}[t]
    \centering
    \includegraphics[width=\textwidth]{abstract_syntax_mm_drawio.drawio.pdf}
    \caption{Abstract syntax metamodel of our DSL.}
    \label{fig:metamodel}
\end{figure*}


%The concepts represented in the metamodel cover the aspects of previous work~\cite{DBLP:conf/icse/IzquierdoC15} and extend them to include the role of agents in the governance policies.
Our proposal is structured around key conceptual areas:

\smallskip
\noindent\textbf{Participants: Embracing Diversity and Agent Characteristics.}
A vital aspect of our DSL is the representation of \texttt{Participant}. %This is crucial for establishing governance that can adapt to the varied nature of team members.
The metamodel distinguishes between human participants (potentially detailed with \texttt{Profile} information like gender or role-specific attributes to support fairness or weighted influence) and \texttt{Agent}s. 
The influence of each participant in the decision-making process can be adjusted through the \texttt{vote\_value} attribute, which can be set to favor certain individuals or groups.
For \texttt{Agent}s, we describe further attributes critical for governance, such as \texttt{autonomy\_level}, \texttt{explainability} and \texttt{confidence}. %, as discussed in Section~\ref{sec:motivation}. 
%These attributes allow policies to be aware of the agent's declared capabilities and reliability, thus adapting the impact of the agents in the decision-making. 
This allows expressing complex policies, such as giving a high-autonomy agent more freedom, while decisions from agents with low explainability might be used only as support.
Participants can also be grouped by \texttt{Role}, while also being flexible to allow certain individuals to act as part of a role in specific policies (see \texttt{hasRole}).\looseness-1

\smallskip
\noindent\textbf{Policies: Defining Flexible and Composable Decision-Making.}
The DSL provides a rich set of \texttt{Policy} constructs to capture diverse decision-making processes essential for adaptive governance.
It supports various policy resolution strategies, including \texttt{VotingPolicy} (e.g., majority and qualified majority based on a \texttt{ratio}), \texttt{ConsensusPolicy} (including \texttt{LazyConsensus}~\cite{LAZYCONSENSUSWEBSITE} as a practical mechanism for efficient agreement), and \texttt{LeaderDrivenPolicy} (with optional fallback mechanisms). 
% This variety allows organizations to specify their existing practices or design new ones suitable for human-agent interactions.
Policies specify a \texttt{DecisionType}, such as \texttt{BooleanDecision} (for accept/reject scenarios like pull request approvals) or \texttt{CandidateChoice} (for selecting one out of multiple options, like electing a leader).
Simple policies can be combined into \texttt{ComposedPolicy} structures (e.g., requiring sequential evaluation or all sub-policies to pass). 
This enables the modular construction of complex governance scenarios from simpler, reusable components, thus enhancing clarity and maintainability.

\smallskip
\noindent\textbf{Scope and Conditions: Contextualizing Policy Application.}
Governance policies are usually applied in a particular project or granular component of it. Our DSL allows their application to be precisely contextualized.
Policies are applied within a defined \texttt{Scope}, such as a \texttt{Project}, an \texttt{Activity} (e.g., development, testing), or a specific \texttt{Task} (e.g., a patch). 
This hierarchical scoping enables granular control, allowing different governance rules for different project areas or activities.
Furthermore, policies can be described with \texttt{Condition}s (e.g., setting a deadline or a minimum required participants). %, or the exclusion of some participants to prevent conflicts of interest). 
These conditions can also act as prerequisites (\texttt{pre}) or checks following execution (\texttt{post}). 



% Mainly, the metamodel is aimed at representing four main aspects: 
% (1) scope (i.e., where the policy is applied);
% (2) participants (i.e., who is involved in the policy), 
% (3) policies (i.e., what the policy is about), and 
% (4) conditions (i.e., when it should be applied). 
% Next we describe the main metaclasses of the metamodel according to these four aspects.

% \smallskip
% \noindent \textbf{Scope.}
% To represent the first aspect, we propose the \texttt{Scope} generalization.
% % We define three types of scopes so that there is a greater granularity at where the policy is applied: 
% We envisioned a three-layered structure consisting of the:
% (1) \texttt{Project} metaclass representing the whole system; % "system" to not say again project
% (2) \texttt{Activity} metaclass representing a subset of a project (e.g., development activities);
% and (3) \texttt{Task} metaclass representing the atomic part of a project (e.g., a pull request in the context of OSS projects).
% This structure allows the extension of the metamodel to include other scopes in the future. % Maybe mention this in a more general way (since we also do it in the other blocks)
% Each metaclass represents one level in the atomicity of the scope, thus they can be associated with each other.
% We also defined a \texttt{StatusEnum} enumeration to represent the state, which can be accepted, meaning the policy instance itself is accepted for application; partial, meaning the process is ongoing; or completed, meaning the process is finished.
% This attribute is optional, but it can be useful to define the status of the scope in the decision engine.

% \smallskip
% \noindent \textbf{Participants.}
% We define the \texttt{Participant} abstract metaclass, which is a generalization of the \texttt{Role} and \texttt{Individual} metaclasses.
% The former is intended to represent a group of individuals that share a common role in the project (e.g., developers, maintainers, etc.), while the latter is intended to represent a specific individual (e.g., an owner).
% To ensure fair representation we define an attribute to favor contributions from certain individuals in a decision (see \texttt{vote\_value}). 
% As discussed in Section~\ref{sec:motivation}, we can further identify individuals by their profile, based on their gender or other characteristics.
% We provided a subset of attributes to represent the profile of an individual (see \texttt{Profile} metaclass), which can be derived from related work~\cite{DBLP:journals/corr/abs-2503-05470}.
% These attributes are optional, but they can be used to calculate the vote value of an individual and favor certain minorities.
% To represent agents (see \texttt{Agent} metaclass) in the project, we extend the \texttt{Individual} metaclass.
% We defined a set of attributes aforementioned discussed (see Section~\ref{sec:motivation}).
% These quantitative indicators (i.e., \texttt{autonomy\_level}, \texttt{explainability}, and \texttt{confidence}) can be used in the  engine in order to assess properly the participation of agents in the decision-making. 
% Additionally, we defined a ternary association between the \texttt{Individual}, the \texttt{Scope} and the \texttt{Role} metaclasses (see \texttt{HasRole} metaclass) to represent the particular role an individual might have on a specific scope (e.g., a user might be contributor for development, but reviewer for documentation).

% \adem{might be a bit weird to reference it to ``considerations'', but I refer to ``the metamodel is composed of four main considerations''. We could use subsections for each of them, and present their syntax with examples.}
% \smallskip
% \noindent \textbf{Policies.}
% To define policies, we define the policy hierarchy, which allows defining two kinds of policies: single policies and composed policies, which consist of a set of policies (either single or composed).
% Depending on the method apply for resolution, the policy can be based on:
% (1) voting, where each participant selects one (or more) candidates among different possibilities; % where we cover two scenarios, majority (by default is evaluated above 50\% but can be modified with \texttt{ratio} attribute) based on the number of participants involved (see \texttt{Majority} metaclass) or majority based on the number of participants that voted (see \texttt{AbsoluteMajority} metaclass);
% (2) consensus, which is focused on a cooperative dynamic where all participants work together to make the best possible decision for the group; % we bring support for lazy consensusCITE (see \texttt{LazyConsensus} metaclass);
% or (3) leader-driven, where one participant (usually the owner) takes the decision. % (note that we can define a fallback policy, see \texttt{default} \textbf{attribute}).

% In our proposal, we generalized the evaluation of voting policies in two kinds: (1) by the number of participants involved in the voting process (see \texttt{AbsoluteMajority} metaclass) or (2) by the number of participants that voted (see \texttt{Majority} metaclass).
% In a nutshell, the decision is made by assessing whether the most voted option surpasses a certain threshold based on the number of potential voters or the number of actual voters.
% For instance, in a voting policy with 10 participants, with a threshold of 50\%, if 4 of them vote for approval and 2 for rejection, the voting policy will be approved if we defined a majority policy, but it will not be approved if we follow an absolute majority policy.
% The threshold by default is set to 50\%, but it can be modified with the \texttt{ratio} attribute.

% Consensus policies are based on the idea that all participants should agree on a decision before it is made~\cite{butler1991conflict}. \adem{This book is not from a conference/journal, can we add it?}
% This is particularly useful in situations where the decision has a significant impact on the project or the community.
% We provided support for lazy consensus~\footnote{\url{https://www.apache.org/foundation/glossary.html#LazyConsensus}}, which is a form of consensus where a decision is made if no one objects.
% It is, as mentioned in the \textsc{Apache} Foundation, one of the most preferred consensus methods.

% Leader-driven policies represent the case where a participant, playing the role of a leader, makes the decision for the group.
% Because of this decision-making process relies on a single individual, the leader can define a fallback policy (see \texttt{default} attribute) in case the decision is not made by the deadline.

% In order to express what the policy are selecting, we defined a set of decision types (see \texttt{DecisionType} hierarchy).
% A policy must have one decision type.
% Policies looking for an acceptance or rejection of a proposal can be represented as a \texttt{BooleanDecision} (e.g., a pull request review).
% Policies looking for a selection among a set of options can be represented as a \texttt{CandidateChoice} (e.g., who will become the new leader of a community).
% Thus, the set of options can be represented as a list of elements of our metamodel (see \texttt{ElementList} metaclass) or as a set of strings (see \texttt{StringList} metaclass).

% These policies can be combined to create more complex policies.
% Composed policies (see \texttt{ComposedPolicy} metaclass) are defined by a set of policies (either single or composed).
% The evaluation of the composed policy is based on the evaluation of its sub-policies.
% For this, we propose a set of boolean attributes to define the relationship between the sub-policies: 
% (1) \texttt{sequential}, sub-policies are evaluated by order of declaration;
% (2) \texttt{require\_all}, all policies must be satisfied;
% and (3) \texttt{carry\_over}, the outcome of the first policy is used as input for the next one (this requires \texttt{sequential} attribute to be true).

% \smallskip
% \noindent \textbf{Conditions.}
% we defined the condition hierarchy (see \texttt{Condition} metaclass) to express the determinants of the policy (e.g., a policy to be resolved before two weeks) and the policy hierarchy (\texttt{Policy} metaclass) to define the method in which to resolve the policy (e.g., absolute majority).
% The condition hierarchy is composed of several conditionals, and structural properties that define the policy (see, for instance, \texttt{ParticipantExclusion} metaclass).
% Conditionals are evaluated before or after the policy is executed (see \texttt{pre} and \texttt{post} attributes).
% If they are fulfilled, the engine calculates the outcome.
% For instance, in a policy that requires three participants (see \texttt{MinimumParticipant} metaclass), it will not be evaluated until this condition is fulfilled (which means that by the end of the deadline, if exists, the policy will be rejected).
% Structural properties are considerations that will be taken in the evaluation of the policy.
% For instance, in a voting policy, we can define an exclusion of certain participants (see \texttt{ParticipantExclusion}) to avoid biases in their approval (e.g., excluding the one proposing the change).
% Similarly, a policy can have a deadline (see \texttt{Deadline} metaclass), defined by a fixed date or an offset (e.g., two weeks).

\subsection{Concrete Syntax}
\label{sec:our-dsl:concrete-syntax}

As concrete syntax, we use a textual notation following a typical block-based structure (see Listing~\ref{lst:example}).
The syntax is therefore formally defined by a grammar. % that specifies the valid language constructs and their composition rules.
Each instance of the metaclass is textually represented by its keyword and a block that contains the properties of the instance.
Containment references are represented as nested blocks while non-containment references use an identifier to refer to the target element.\looseness-1

% To define an instance of a governance policy, we need to define the scope of the policy, the participants involved, the conditions that must be fulfilled, and the method in which to resolve the policy.
% Listing~\ref{lst:example} shows a simple governance policy evaluated following the majority method as a \texttt{BooleanDecision} requiring 40\% positive votes, and it has a deadline of ten days and excludes one participant.

% \adem{If space is needed, we can shorten/remove this paragraph explaining the listing.}
% As we see, we first define the scopes and their relationships.
% % An alternative would be to just define one scope without their containments, but this would lack information for the decision engine. 
% % \adem{This sentence is to say, we can define a Task, but how will we know to where it pertains so that we can execute the engine? Maybe, since we later focus on GH, this is not that relevant to say}
% Next, we declare the possible participants, in our case a \texttt{Maintainer} role and two individuals, \texttt{Joe} and \texttt{Mike} with their attributes. 
% We can also define further the profile of the individuals. 
% % We define the gender of \texttt{Joe} as \texttt{male}.
% In the third block we define the policy object.
% The first keyword is the policy type, in this case it is a \texttt{MajorityPolicy}, followed by its name.
% Inside the policy block, we declare the relationships to the scope and the participants as a list.
% Note that inside the participant list we can define the role an individual might have inside a policy with the \texttt{as} keyword followed by the role defined before.
% A policy might have certain conditions, these are defined with the \texttt{Conditions} keyword.
% % In our example, we define the deadline of ten days, we exclude \texttt{Mike} of the voting count with the \texttt{ParticipantExclusion} keyword, and we require at least three voters, defined with the \texttt{MinParticipants} keyword.
% At last, any parameter the policy might have, we declare it in the \texttt{Parameters} attribute, followed by the set of parameters, which in our case is the \texttt{ratio} set at 40\%.

\begin{lstlisting}[caption=Simple governance policy evaluated following a majority strategy as a \texttt{BooleanDecision} requiring 40\% positive votes with a deadline of ten days., 
                    label=lst:example, 
                    basicstyle=\ttfamily\fontsize{8pt}{8.5pt}\selectfont,
                    breaklines=true,
                    float=t,
                    captionpos=b]
Scopes: 
    Project myProject {
        activities : myActivity {
            tasks : myTask
        }
    }
Participants:
    Profiles : 
        profile1 {
            gender : male
            race : hispanic
        }
    Roles : Maintainer
    Individuals : 
        Joe {
            vote value : 0.7
            profile : profile1
            role: Maintainer
        }, 
        (Agent) Mike {
            confidence : 0.8
            role: Maintainer
        },
        Paul {
            role: Maintainer
        }
MajorityPolicy TestPolicy {
    Scope: myTask
    DecisionType as BooleanDecision
    Participant list : Maintainer
    Conditions:
        Deadline : 10 days
        ParticipantExclusion : Paul
    Parameters:
        ratio : 0.4
}
\end{lstlisting}
\subsection{Operational Semantics: Decision engine}
The operational semantics take the form of a decision engine, enforcing the policies on ongoing collaborations.
The engine receives platform data to update the decision-making process state and take the necessary actions.
For the first implementation of our approach, we rely on \textsc{GitHub} as the collaboration platform.
% , using its webhook mechanism and REST API to monitor interactions and enact decisions.

The decision engine operates through a continuous cycle of state management and policy enforcement. 
When a collaboration event occurs (e.g., a pull request is created), the engine captures the event data and creates an internal representation. %, including its scope, rationale, and associated metadata. 
This state management enables the engine to track multiple concurrent decision-making processes.\looseness-1

% When receiving a new collaboration, the engine searches for the applicable policy based on the collaboration scope.
For new collaborations, the engine identifies applicable policies based on scope.
Once found, the engine creates an internal deadline check based on the policy to resolve the 
%decision-making 
process.
% Then, when a voting action is performed
When voting occurs (e.g., a pull request review), 
eligible users' votes are registered in the collaboration's ballot box.
% if the user is allowed to vote (according to an ongoing policy), the engine registers the vote, adding it to the associated ballot box of the collaboration.
Each vote is captured with its rationale, timestamp, and association to the voter. %, building a complete audit trail of the decision-making process.
% Finally, when the deadline is reached, 
At deadline, the engine analyzes the votes 
% based on the policy and conditions 
and enacts the decision (e.g., merge the pull request).
% In the context of composed policies, their resolution is based on the resolution of their phases.
% The votes required by the phases happen sequentially or in parallel based on the \texttt{sequential} attribute, the votes in a ballot box are transferred to the following one if \texttt{carry\_over} is true, and the \texttt{require\_all} attribute is used to make the composition of the phases' decisions (i.e., conjunction if true, disjunction if false).
% In the case of composed policies, the decisions of the phases are resolved 
Composed policies resolve phases
sequentially or in parallel, 
%and composed based on the requirement of all decisions approval or not (i.e., conjunction or disjunction).
combining decisions via conjunction or disjunction.
In addition, 
% the votes in a ballot box can be carried over to the next phase ballot box.\looseness-1
votes can be carried over between phases.


%%%%
%% Description based on classes
% At the core of the engine, collaboration models represent the state of all previous and ongoing collaborations.
% % This state is stored following the schema depicted in Figure~\ref{fig:CollabMetamodel}. 
% % The \texttt{Interaction} class aggregates the set of collaborations happening, the individuals involved and the decisions taken.
% \texttt{Collaboration} represents the managed collaborations, which are defined using an identifier, a rationale, a scope, a set of ballot boxes, and an association to the \texttt{Decision} for this collaboration.
% \texttt{Scope} represents the scope in which the collaboration was created.
% The ballot boxes represent the mapping between applicable policies and the set of votes associated with this policy.
% This mechanism allows managing composed policies that require multiple voting resolutions.
% These votes are represented using \texttt{Vote} metaclass, characterized by an agreement value, a natural language rationale, and a timestamp.
% Votes are associated with an \texttt{Individual} (the voter), one or more ballot boxes, and the resulting decisions.
% \texttt{Decision} models decisions, registering the result and timestamp of the decision process, the policy applied, and references to the applying votes.

% The collaboration model is updated based on the information provided by the collaboration platform.
% First, when a collaboration begins %on the platform 
% (i.e., pull request), the decision engine creates the corresponding \texttt{Collaboration} instance using the data provided.
% In addition, the engine searches for the applicable policy based on the collaboration scope.
% When the appropriate policy is found, the engine creates an internal deadline check based on the policy to resolve the decision-making process.
% Then, when a voting action is performed (i.e., pull request review), if the user is allowed to vote (according to the policy), the vote object is created and added to the appropriate ballot boxes of the collaboration.
% Finally, when the deadline is reached, the engine will analyze the votes based on the policy and conditions and create the associated decision object and enact the decision (i.e., merge the pull request).

% In the context of phased policies, their resolution is based on the resolution of their phases.
% The votes required by the phases happen sequentially or in parallel based on the \texttt{sequential} attribute, the votes in a ballot box are transferred to the following one if \texttt{carry\_over} is true, and the \texttt{require\_all} attribute is used to make the composition of the phases' decisions (i.e., conjunction if true, disjunction if false).

\section{Proof of Concept}
\label{sec:proof-of-concept}
We implemented the metamodel as a set of Python classes, while the concrete syntax was implemented through a grammar definition using ANTLR~\cite{ANTLRWEBSITE}, %a parser generator that enables the creation of 
% a tool for creating language recognizers, interpreters, and translators from grammars. 
a parser generator tool.
We developed a transformation that converts the Abstract Syntax Tree (AST) generated by ANTLR into instances of our metamodel, which serve as input for the decision engine.
% to the governance models as instances of the metamodel to be then used as input for the decision engine.
The complete implementation, including the grammar and decision engine, are available in the tool repository.\footnote{\url{https://doi.org/10.5281/zenodo.15856633}}
We also translated several governance policies identified in our background analysis (see Section~\ref{sec:motivation:governance-oss}) into our DSL syntax to demonstrate its expressiveness and practical applicability. 


% \section{Enforcing governance on GitHub}
% \label{sec:engine}
% \input{ase-engine}

% \section{Validation}
% \label{sec:validation}
% \input{ase-validation}

% \section{Related Work}
% \label{sec:related}
% \input{ase-related}

\section{Roadmap}
\label{sec:discussion}
The proposed DSL is just a first step towards a more diverse,  
transparent and agentic collaboration in software development. 
Open challenges and future directions are discussed below.

\smallskip
\noindent\textbf{Usability improvement via new syntaxes.}
To facilitate the adoption of our DSL by less technical users, we envision offering alternative UIs, like chatbots, to enhance usability.
%  to define the policies, enhancing usability through intuitive UIs or chatbot interfaces. 
This will be especially useful when introducing our DSL to manage governance in other, non-software, domains.
% Furthermore, we could leverage some technologies, such as BESSER Agentic Framework\footnote{\url{https://github.com/BESSER-PEARL/BESSER-Agentic-Framework}}, to create a chatbot interface that guides users through the policy definition process.
%This would broaden accessibility beyond technical users, fostering more inclusive governance design.

\smallskip
\noindent\textbf{Longitudinal study.}
To better understand if explicit governance rules help to increase 
%the number of contributions in an OSS project,
OSS contributions,
% (and, if so, what types of policies are more helpful to that), 
we plan a longitudinal study where we could compare the evolution of contributions in projects before and after they committed to an explicit governance model. 
This same study could be useful to see whether projects are sticking to such governance model or they just ignore it even after committing to it. 

%ly, we intend to investigate usability aspects and gather feedback from practitioners to refine the language and its supporting tools.

\smallskip
\noindent\textbf{Improved impact of Multi-agent systems in software development.}
We plan to test whether our governance layer enables more complex multi-agent systems to collaborate, in a controlled way, to solve more ambitious software engineering tasks. 
By leveraging SWE-bench~\cite{DBLP:conf/iclr/JimenezYWYPPN24}, we can systematically evaluate how governance policies affect agent collaboration, task completion rates, and adherence to defined constraints.
% This evaluation will provide insights into how formal governance structures impact autonomous agent behavior and overall system performance in software engineering contexts.

\smallskip
\noindent\textbf{Participant identification and bias mitigation.}
A problem some projects face is identifying the right people to play certain roles. 
Given a governance model, it would be possible to analyze past project data to propose candidates to fill the roles required in the governance model. 
% These candidates could then be validated by the community. 
A real-time monitoring component should also be helpful to proactively change the candidates to continuously improve the diversity of the collaborations. Indeed, a key element of this (semi)automatic participant analysis could be devoted to make sure the participants list is not biased against certain communities and that it respects the profile requirements. 
%This would enable automated role adjustments, proactive governance violation detection, and adaptive policy enforcement incorporating positive and negative discrimination strategies to maintain fair representation.

%\smallskip
%\noindent\textbf{Bias mitigation and diversity enhancement.}
%Finally, we plan to introduce mechanisms that actively promote diversity in governance decisions and team composition.
%This includes automatic role identification in governance processes and implementing diversity metrics that can be continuously monitored and enforced, thus maintaining inclusive participation across different profile dimensions.

\smallskip
\noindent\textbf{Governance beyond software.}
We believe other domains could also benefit from an explicit and precise definition of their governance policies based on our DSL. 
% This broader applicability underscores the potential impact of a generalized governance framework for mixed-initiative systems and is a key direction for future exploration.
For instance, making explicit the governance systems of organizations (such as NGOs) would enable systematic comparison of their decision-making structures and their assessment of inclusivity. % across different organizational contexts.

\section{Conclusion}
\label{sec:conclusion}
We have presented a vision of a DSL for the explicit definition and automation of governance policies in software projects, with a particular focus on supporting diversity among participants, including both human contributors and AI-powered agents. 
Our DSL extends previous work by enabling the specification of participant profiles and agent attributes, thus facilitating more equitable and transparent decision-making processes. 
% Through the analysis and translation of governance policies from real-world projects, we have demonstrated the expressivity and adaptability of our approach, as well as its applicability to contemporary collaborative software development environments.
This is just a first step in this direction. As future work, we plan to work on the aforementioned roadmap.



\bibliographystyle{ieeetr}
\bibliography{ase-main}

\end{document}
